\documentclass[11pt,a4paper]{article}

\usepackage[utf8]{inputenc}
\usepackage[T1]{fontenc}
\usepackage[margin=2.5cm]{geometry}
\usepackage{graphicx}
\usepackage{booktabs}
\usepackage{array}
\newcolumntype{L}[1]{>{\raggedright\arraybackslash}p{#1}}
\usepackage{longtable}
\usepackage{listings}
\usepackage{xcolor}
\usepackage{enumitem}
\usepackage[hidelinks]{hyperref}

% --- light code styling, nothing fancy -------------------------------------
\definecolor{codebg}{gray}{0.96}
\definecolor{codecomment}{gray}{0.45}
\lstset{
  basicstyle=\ttfamily\footnotesize,
  backgroundcolor=\color{codebg},
  commentstyle=\color{codecomment}\itshape,
  keywordstyle=\bfseries,
  breaklines=true,
  frame=none,
  showstringspaces=false,
  columns=fullflexible,
  xleftmargin=8pt,
  aboveskip=8pt,
  belowskip=8pt,
}
\lstdefinestyle{py}{language=Python}
\lstdefinestyle{sh}{language=bash}

\newcommand{\file}[1]{\texttt{#1}}
\setlist{itemsep=2pt, parsep=0pt, topsep=4pt}

\title{Circularity Dashboard\\
\large A France--Brazil comparison of Extended Producer Responsibility\\[4pt]
\normalsize Code and data documentation}
\author{Arthur Briens \\ \small USP--PSL Institute for Circular Economy}
\date{\today}

\begin{document}
\maketitle

\begin{abstract}
\noindent
This document explains how the circularity dashboard is built: what each file
does, why it is written the way it is, and which traps in the underlying data
the code is deliberately defending against. It is written to be read by someone
who has not seen the code before, and to be useful again in six months when the
reasoning behind a particular line has been forgotten. Sections~\ref{sec:run}
and~\ref{sec:model} are the ones to read first; the rest can be consulted as
reference.
\end{abstract}

\tableofcontents
\newpage

% ===========================================================================
\section{What this project is}

The dashboard compares how France and Brazil regulate the circular economy. It
has three substantive pages:

\begin{itemize}
  \item \textbf{Targets} --- what each country has committed to, in law, and how
        far it has got.
  \item \textbf{EPR market structure} --- how concentrated the French compliance
        market is: monopoly, duopoly or oligopoly, filière by filière.
  \item \textbf{France, EPR filières} --- tonnages placed on the market and
        treated, by filière, eco-organisme and region.
\end{itemize}

It is a Streamlit application written in Python, reading ADEME open data for the
French side and a hand-assembled, sourced CSV for the targets comparison.

\subsection{The one idea behind the whole thing}

France and Brazil do not set targets on the same footing, and the easiest way to
produce a misleading comparison is to put two percentages on one axis when their
denominators differ. France measures WEEE collection as 65\,\% of a rolling
three-year average of equipment placed on the market. Brazil measures it as
17\,\% of product weight against a fixed 2018 base year. Both are ``a
percentage of what was put on the market'', both can be plotted as a bar, and
they are not the same quantity.

Almost every design decision in this codebase follows from taking that problem
seriously. The data model carries a \file{basis} field (the denominator, written
out in words) alongside every number, and the charts refuse to place two
countries on a shared axis unless their bases actually match.

% ===========================================================================
\section{Getting it running}\label{sec:run}

\subsection{Locally}

From the project root:

\begin{lstlisting}[style=sh]
pip install -r requirements.txt
streamlit run output/dashboard/app.py
\end{lstlisting}

The app opens in a browser. The three pages appear in the left-hand navigation,
which Streamlit builds automatically from the files in
\file{output/dashboard/pages/}. The numeric prefix in \file{2\_targets.py} sets
the ordering; it is not part of the page name.

\subsection{Deployed}

The app is deployed on Streamlit Community Cloud from a private GitHub
repository. The entry point is \file{output/dashboard/app.py}. Cloud rebuilds
from a fresh clone on every push.

\subsection{Where the data comes from}\label{sec:sources}

The three ADEME extracts are deliberately \emph{not} committed to git; the
targets CSV \emph{is}. The line drawn is \textbf{third-party bulk extracts out,
our own research in} --- the ADEME files total roughly 4.5\,MB and are
redownloadable at any time, whereas \file{targets\_fr\_br.csv} is 84\,KB of
hand-assembled, individually sourced rows that the targets page needs in order
to be reproducible from a clone.

Either way, each file is resolved in two steps:

\begin{enumerate}
  \item \file{data/raw/<filename>} --- the local working copy, used in
        development;
  \item a URL held in \file{st.secrets["data"]} --- used by the deployed app.
\end{enumerate}

A fresh clone contains the targets CSV but none of the ADEME extracts, so those
fall through to the secrets URL. Configure them on Streamlit Cloud under
\textbf{Manage app $\rightarrow$ Settings $\rightarrow$ Secrets}:

\begin{lstlisting}
[data]
msm_url     = "https://.../REP.csv"
trt_url     = "https://.../traitements_REP.csv"
ref_url     = "https://.../Referentiels_MSM.xlsx"
targets_url = "https://.../targets_fr_br.csv"   # optional: overrides the
                                                # committed copy
\end{lstlisting}

Locally the same block goes in \file{.streamlit/secrets.toml}, which is
gitignored. \file{secrets.toml.example} at the repository root documents the
shape and is safe to commit.

\paragraph{These URLs are credentials.} A pre-signed S3 link or a share link
with an embedded token grants read access to whoever holds it. They must never
enter git, a log line, or an error message. \file{lib/data.py} contains a
\file{\_scrub()} helper whose only job is to strip URLs out of exception text
before it is displayed, because an error message is the easy way to leak one.

% ===========================================================================
\section{The mental model}\label{sec:model}

Two ideas explain most of the code. Neither is obvious from reading it.

\subsection{Streamlit re-runs the whole script}

Streamlit re-executes the \emph{entire page script, top to bottom, on every
interaction}. Move a slider and the file runs again from line 1. There is no
callback system and no component tree that persists between interactions.

Three consequences follow, and they explain a lot:

\begin{itemize}
  \item \textbf{Widgets are just values.} \file{dark = st.toggle(...)} returns a
        plain boolean for this run. There is no state plumbing.
  \item \textbf{Expensive work must be cached}, or a 2.6\,MB CSV is re-read on
        every click. Hence \file{@st.cache\_data} on every loader.
  \item \textbf{Declaration order is render order.} A control created after
        \file{st.columns()} appears \emph{below} the charts it drives. This
        caused a real bug: the ``filière in focus'' selector rendered under the
        charts it controlled until it was moved above the column split.
\end{itemize}

\subsection{The visual layer is three tiers that do not overlap}

Styling a Streamlit app reaches different places depending on which lever you
pull, and it is worth knowing which is which.

\begin{longtable}{@{}L{3.9cm}L{5.0cm}L{5.4cm}@{}}
\toprule
\textbf{Tier} & \textbf{Reaches} & \textbf{Limitation} \\
\midrule
\endhead
\file{.streamlit/}\newline\file{config.toml} &
Streamlit's own widgets: sidebar, sliders, multiselects, buttons &
Read once at server start. \emph{Static} --- cannot respond to anything in the
script \\
\addlinespace
Plotly template
(\file{lib/theme.py}) &
Axis colours, gridlines, fonts, hover styling, legend placement, the
categorical colour order &
Only applies to figures, and Streamlit will override parts of it unless told
not to (see below) \\
\addlinespace
Injected CSS
(\file{theme.page\_css}) &
Metric cards, the hero number, section labels, chart card backgrounds, text
colour in dark mode &
Targets Streamlit's internal \file{data-testid} attributes, which can change
between versions. The fragile tier --- kept to a handful of rules \\
\bottomrule
\end{longtable}

\paragraph{The seam.} Because tier~1 is static and tiers~2--3 are dynamic, the
in-app dark-mode toggle changes the charts and cards but \emph{not} Streamlit's
sidebar and widgets. That is a genuine half-measure. Making it complete would
mean setting \file{base = "dark"} in \file{config.toml}, which is an app-wide,
restart-level setting rather than a per-session toggle.

\paragraph{Two Plotly traps, both discovered the hard way.}
\begin{enumerate}
  \item \file{st.plotly\_chart} defaults to \file{theme="streamlit"}, which
        overrides your registered template. Every call in this codebase passes
        \file{theme=None}.
  \item Even then, Streamlit merges \file{config.toml}'s colours over the
        figure's \emph{template} values, because template values count as
        defaults. Only an explicit value on the figure's own layout wins that
        merge. So every chart helper sets
        \file{paper\_bgcolor} and \file{plot\_bgcolor} explicitly. Without this,
        dark mode rendered every plot area white under a dark page.
\end{enumerate}

% ===========================================================================
\section{Repository map}

\begin{longtable}{@{}L{5.6cm}rL{7.2cm}@{}}
\toprule
\textbf{Path} & \textbf{Lines} & \textbf{Role} \\
\midrule
\endhead
\file{output/dashboard/app.py} & 95 & Landing page. Project description and
links to the three pages \\
\file{lib/\_\_init\_\_.py} & 1 & Makes \file{lib} a package \\
\file{lib/data.py} & 514 & Source resolution, ADEME loaders, French transforms \\
\file{lib/market.py} & 314 & Market-structure measures (monopoly/oligopoly) \\
\file{lib/targets.py} & 549 & Targets schema, placeholder set, loader, pairing \\
\file{lib/theme.py} & 206 & Palette, Plotly templates, page CSS \\
\file{pages/2\_targets.py} & 492 & France vs Brazil target comparison \\
\file{pages/3\_EPR.py} & 568 & EPR market structure \\
\file{pages/4\_France\_EPR.py} & 446 & France filière detail \\
\midrule
\file{data/raw/REP.csv} & 1.8\,MB & Put-on-market declarations, 2017--2024 \\
\file{data/raw/traitements\_REP.csv} & 2.6\,MB & Treated-waste declarations \\
\file{data/raw/Referentiels\_MSM.xlsx} & 76\,KB & Actor and filière lookups \\
\file{data/raw/targets\_fr\_br.csv} & 81\,KB & 194 researched targets \\
\midrule
\file{docs/targets\_csv\_schema.md} & --- & Specification for the targets CSV \\
\file{requirements.txt} & --- & pandas, openpyxl, streamlit, plotly \\
\file{secrets.toml.example} & --- & Template for the data-source URLs \\
\file{.streamlit/config.toml} & --- & Streamlit theme (accent colour, fonts) \\
\file{code/*.py} & --- & Earlier exploratory analysis, not used by the app \\
\bottomrule
\end{longtable}

The split is deliberate: \file{lib/} contains everything that touches pandas and
nothing that touches presentation; \file{pages/} contains composition and almost
no data logic. A practical benefit is that you can debug a chart without
launching Streamlit at all:

\begin{lstlisting}[style=py]
from lib import data as D, market as M
panel = M.actor_panel(D.load_msm(), D.load_actors())
M.structure(panel).query("filiere == 'EEE'")
\end{lstlisting}

\newpage
% ===========================================================================
\section{The library layer}

\subsection{\file{lib/data.py} --- getting the ADEME data in}

Three responsibilities: find the files, load them, reshape them.

\subsubsection*{Finding the files}

A small \file{Source} dataclass names each required file, its expected local
filename, and its key in the secrets table. \file{\_open(key)} returns either a
\file{Path} or an in-memory buffer, and pandas accepts both, so the loaders do
not care which they got.

\file{require\_data()} is called at the top of every page that needs data. It
does two things, both about failing legibly:

\begin{enumerate}
  \item checks that each source is configured, and if not, prints a message
        naming the missing file and both ways to fix it;
  \item \emph{primes} the loaders inside a \file{try}, so that a download that
        times out or returns 403 is reported here rather than surfacing deep
        inside a chart.
\end{enumerate}

This matters because Streamlit Community Cloud \textbf{redacts exception text}
to avoid leaking data. An unguarded pandas error reaches the user as ``the
original error message is redacted'' with no indication of which file is
missing --- which is exactly the failure this project hit in deployment.

\file{require\_data()} is deliberately \emph{not} cached: \file{st.stop()} raises
a control-flow exception, and caching a function that raises it would poison the
cache entry.

\subsubsection*{Path resolution}

\begin{lstlisting}[style=py]
ROOT = Path(__file__).resolve().parents[3]   # lib -> dashboard -> output -> repo
RAW  = ROOT / "data" / "raw"
\end{lstlisting}

Paths resolve from the file, not the working directory. A relative
\file{"data/raw/REP.csv"} only works when the app happens to be launched from
the project root, which is not guaranteed on a host.

\subsubsection*{Reshaping, and the traps it defends against}

The transform functions each return a tidy frame shaped for exactly one chart.
Three defensive details are worth knowing, because each was a real bug:

\begin{itemize}
  \item \textbf{\file{type\_ton} is null for most filières.} Any \file{groupby}
        or merge on it silently drops EEE, TLC and EA. It is filled with
        \file{"N/A"} before it is ever used as a key.
  \item \textbf{\file{pays\_site\_trt} is null for all EMPAP rows.} Treating
        null as ``not France'' rendered EMPAP as 100\,\% exported, which is
        nonsense. The field is now three-state: France, abroad, or
        \emph{not declared}.
  \item \textbf{\file{reg\_site\_trt} uses pre-2016 region codes} --- 22
        régions, not the current 13. \file{LEGACY\_REGION} maps them across.
\end{itemize}

\subsection{\file{lib/theme.py} --- one visual system}

Pure parameters, no data knowledge. Two palette dictionaries (light and dark), a
Plotly template factory, and a CSS string builder.

The categorical colour order is validated for colour-vision deficiency in both
modes. \textbf{The order is the safety mechanism, not decoration.} An attempt to
re-order the palette so that two particular hues were not adjacent failed
validation badly in dark mode, and was reverted.

Three colour roles, and mixing them up is the most common way a chart goes
wrong:

\begin{itemize}
  \item \textbf{Categorical} --- identity, no order. Countries, eco-organismes.
  \item \textbf{Ordinal} --- ranked categories. The waste hierarchy; the
        monopoly-to-oligopoly scale. Darkest is ``most concentrated''.
  \item \textbf{Sequential} --- continuous magnitude. Heatmaps only.
\end{itemize}

Status colours (good/warning/critical) are reserved and never reused as a
series colour. This matters in practice: Plotly assigns marker colour from the
colourway if you do not specify one, so target markers silently came out green
--- reading as ``good'' next to a chart about missing targets --- until they
were coloured explicitly.

\subsection{\file{lib/market.py} --- monopoly, duopoly, oligopoly}

A REP filière is a regulated market: producers must join an eco-organisme or run
an approved individual system, so the declarants in \file{REP.csv} are that
market's suppliers of compliance. The standard industrial-organisation toolkit
therefore applies.

\paragraph{Two classifications, deliberately.}

\begin{itemize}
  \item \textbf{Structure} counts declarants: monopoly (1), duopoly (2), tight
        oligopoly (3--4), broad oligopoly (5+).
  \item \textbf{Effective} looks at what the leader actually holds: sole
        declarant, effective monopoly ($\geq 95\,\%$), dominant leader
        ($\geq 50\,\%$), or contested.
\end{itemize}

They disagree, and the disagreement is the finding. PCHIM has three accredited
declarants --- a ``tight oligopoly'' by any count --- and a leader holding
essentially the whole market. Both thresholds are named constants so they are
easy to argue with.

\paragraph{The coverage trap.} Filières enter the SYDEREP reporting system at
different dates: PNEU from 2017, EEE from 2019, ABJ from 2023. A filière with
zero declarants in a year is \emph{not} a market with no suppliers --- it is a
market not yet reporting. Every function returns only reported filière-years,
and the charts render the rest as blank. Reading a blank as a zero is the
easiest mistake to make with this dataset.

\paragraph{Entry and exit} are measured inside each filière's \emph{own}
reporting window. A body present in the filière's first reported year is the
incumbent at the point the record begins, not an entrant. Measuring against the
dataset's overall window would invent a wave of entries in whichever year each
filière joined the system.

\subsection{\file{lib/targets.py} --- the comparison schema}

The point of this module is the schema, not the numbers.

Each row carries \emph{both} a verbatim \file{metric} (as the instrument words
it) and a normalised \file{metric\_family} (how we choose to group it), plus an
explicit \file{basis} string and a \file{comparable} flag.
\file{metric\_family} is our editorial judgement, not the law's --- it is the
honest place for the comparison to be contestable, and it lives in a
module-level list so it is easy to change.

\paragraph{The pairing rule.} \file{paired()} matches on \file{sector} alone,
deliberately \emph{not} on (sector, metric\_family). France measures packaging
as a recycling rate and Brazil as reverse-logistics recovery, so keying on both
fields matches nothing and the chart comes up empty --- hiding the very
comparison the page exists for. Instead:

\begin{enumerate}
  \item if both countries share a metric family for that sector, restrict to it
        and take the furthest horizon within it;
  \item otherwise fall back to the furthest horizon overall and set
        \file{cross\_family}, so the chart can warn that the two bars measure
        different quantities. That is a stronger warning than a basis mismatch.
\end{enumerate}

\paragraph{A boolean trap worth remembering.} \file{comparable} was originally
parsed with \file{Series.astype(bool)}. If a single row breaks the strict
\file{TRUE}/\file{FALSE} pattern --- a \file{no}, an \file{n/a}, a blank ---
pandas reads the column as text, and \file{astype(bool)} makes \emph{every
non-empty string} \file{True}, including the literal \file{"false"}. Every pair
would have been silently marked as sharing a denominator, which is precisely the
error the page exists to prevent. It now parses text explicitly and treats
anything unrecognised as \file{False}, the safe direction.

% ===========================================================================
\section{The pages}

\subsection{\file{app.py} --- landing page}

Project description, institutional context, and \file{st.page\_link} buttons to
the three pages. No data access, so it always loads even when the data sources
are misconfigured.

\subsection{\file{2\_targets.py} --- France vs Brazil targets}

Two sections, mirroring how the two countries legislate.

\textbf{Section 1, government targets.} France: EU directives transposed
nationally, plus the loi AGEC, plus regional plans. Brazil: the PNRS framing
law, the Planares national plan, and state plans. Two charts:

\begin{itemize}
  \item \emph{When each commitment bites} --- a dumbbell from baseline year to
        target year. This is the only chart on the page with \textbf{no}
        comparability caveat, because time is genuinely shared even when the
        metrics are not. It surfaces something the numbers cannot: France's
        commitments cluster at 2025--2035, Brazil's stretch to 2040.
  \item \emph{Distance still to travel} --- bullet bars where each row is scored
        against \emph{its own} target only, so no cross-country axis comparison
        is implied.
\end{itemize}

\textbf{Section 2, sector-specific targets.} France's per-filière obligations
are market-wide and mandatory; Brazil's \emph{acordos setoriais} bind signatory
firms. That difference is structural and survives any amount of data cleaning.

Markers used in the paired chart:

\begin{itemize}
  \item $\neq$ --- the two countries measure this sector on different
        denominators;
  \item $\triangle$ --- they do not even measure the same quantity, so there is
        no shared metric family to pair on. Read those as ``both regulate this
        sector'', never as ``one is ahead''.
\end{itemize}

\paragraph{Provenance is a first-class feature.} A meter at the top reports what
share of rows carry \file{status = verified}. The \file{status} field has three
values: \file{placeholder} (invented, layout only), \file{sourced} (real,
cited, but not read in the primary legal text) and \file{verified} (read in the
instrument's own text). Only the last counts. The middle value exists because
``placeholder'' would misdescribe a row that is real but second-hand, and
``verified'' would overclaim it.

\subsection{\file{3\_EPR.py} --- EPR market structure}

Reads the put-on-market declarations as an industrial-organisation problem: how
much choice does a producer actually have?

\begin{itemize}
  \item \textbf{A} --- headline counts: single-declarant filières, effective
        monopolies, median leader share, genuinely contested filières.
  \item \textbf{B} --- structure map, filière $\times$ year, on the ordinal ramp
        with monopoly darkest. EEE's staircase reads at a glance: monopoly
        (2019) $\rightarrow$ duopoly (2020) $\rightarrow$ tight oligopoly (2021)
        $\rightarrow$ broad oligopoly (2022).
  \item \textbf{C} --- the chart that does the real work: each filière split
        into leader / 2nd / 3rd / all others, sorted by leader share, with the
        declarant count in brackets. This is where nominal and effective
        structure visibly diverge.
  \item \textbf{D} --- a transition log and an entry/exit dot chart.
  \item \textbf{E} --- one filière in detail: share bands over time, and a
        timeline of which declarants were active when.
  \item \textbf{F} --- operators spanning several filières. Single-filière views
        hide the bodies accredited across many at once.
  \item \textbf{G} --- route to market, from the \file{statut\_prod} field.
\end{itemize}

\subsection{\file{4\_France\_EPR.py} --- filière detail}

Reading order: status $\rightarrow$ trajectory $\rightarrow$ mechanism
$\rightarrow$ territory. Headline tiles, a gap-to-target bullet chart, indexed
trajectories as small multiples, treatment composition on the waste-hierarchy
ramp, eco-organisme shares, domestic versus exported treatment, and a schematic
tile map by region.

Two honest limits are stated on the page itself. Every rate here is
\file{treated mass / tonnage put on market}, a proxy, because the separate
collection dataset is not among the extracts. And the target values in the
gap chart are placeholders held in \file{lib/data.PLACEHOLDER\_TARGETS}.

\paragraph{Why the map is a tile grid, not a choropleth.} Every région gets
equal area, so colour is not biased by territory size and Île-de-France does not
vanish beside Nouvelle-Aquitaine. The value is printed in each tile, so colour
never carries the information alone. It also handles the overseas régions
without the usual inset awkwardness --- relevant, because their collection rates
are markedly lower.

% ===========================================================================
\section{The data files}

\subsection{\file{REP.csv} --- put on market}

29{,}485 rows, 2017--2024, 18 filières. One row per declaring actor per filière
per year. Key columns: \file{annee}, \file{filiere}, \file{acteur},
\file{statut\_prod}, \file{quantite}, \file{tonnage}.

\file{acteur} is a code such as \file{FR000017}; \file{Referentiels\_MSM.xlsx}
maps it to a company name. The declarants are a mix of eco-organismes and
individual systems, and both are kept: an individual system is a genuine
alternative to joining a collective body, so excluding them would overstate
concentration.

\subsection{\file{traitements\_REP.csv} --- treated waste}

34{,}971 rows. \textbf{Semicolon-separated}, unlike \file{REP.csv}. Adds
\file{typ\_trt} (73 treatment codes, including legacy variants),
\file{reg\_site\_trt}, \file{dep\_site\_trt} and \file{pays\_site\_trt}.

Two things to hold on to. First, the location fields refer to where waste was
\textbf{treated}, not collected --- a région with high treated tonnage has
treatment plants, not necessarily good collection. Second, the 73 treatment
codes are mapped to four waste-hierarchy buckets by ordered prefix rules in
\file{lib/data.\_RULES}; anything unmatched shows as \emph{Unclassified} rather
than being folded into a bucket. If that share is large, fix the rules before
trusting the composition chart.

TLC's 2024 treatment is entirely \file{TRI\_INCONNUE} --- sorted, destination
unknown. It renders as a full grey Unclassified bar. That is the data telling
you something real about textile traceability, not a mapping failure.

\subsection{\file{targets\_fr\_br.csv} --- the targets}

194 rows, 19 columns, every row carrying a source URL. 79 government-level and
115 sector-level; 66\,\% marked \file{verified}. This is the one data file
committed to the repository, for the reason given in Section~\ref{sec:sources}. See
Appendix~\ref{sec:schema} for the column specification and
\file{docs/targets\_csv\_schema.md} for the full version.

Notable content: for France, the EU directives (Waste Framework, Landfill,
Packaging including all six per-material targets, Single-Use Plastics, WEEE,
Batteries, ELV, PPWR), the loi AGEC, décret 3R, and the REP \emph{cahiers des
charges} for nine filières. For Brazil, Planares Metas 1--9 with their
2024/2028/2032/2036/2040 trajectories, the staggered \emph{lixão} deadlines, and
the reverse-logistics instruments for plastics, EEE, tyres, lubricants, lamps,
medicines and agrochemicals.

\paragraph{Three caveats carried in the file itself.}
\begin{itemize}
  \item The Planares metas were transcribed from documents reproducing the
        decree's annex, not read in the annex itself. The 2040 endpoints are
        corroborated; the intermediate years are the weakest cells.
  \item Brazil's two national data sources disagree systematically. SINISA puts
        recovery at 1.82\,\% (2023); ABREMA at 8.7\,\% (2024). The gap is
        structural --- ABREMA imputes the tonnage handled by roughly 700{,}000
        informal \emph{catadores}, which SINISA cannot see. Both are carried as
        separate rows rather than picking one.
  \item \textbf{Not one pair shares a denominator.} \file{comparable} is
        \file{FALSE} on all 194 rows. That is a finding, not an omission.
\end{itemize}

% ===========================================================================
\section{Design rules that recur}

Five rules show up across every page. They are worth stating once.

\begin{enumerate}
  \item \textbf{Blank is not zero.} Missing data renders as an empty cell, never
        as a zero bar. Both heatmaps had to be fixed for this: one printed the
        literal \file{NaN}, the other the literal \file{null}.
  \item \textbf{Name the denominator.} Every rate carries a \file{basis} string
        in words. Two percentages may only share an axis when their bases match.
  \item \textbf{Say what is missing.} When a filière drops out of a ranking for
        want of data, the page names it. A row that silently vanishes reads as
        ``not applicable'' rather than ``no data''.
  \item \textbf{Distinguish absence from zero in the data model too.} Brazil has
        \emph{no} battery collection target; the CSV records that explicitly
        rather than leaving the cell blank and letting the reader infer.
  \item \textbf{Sub-1\,\% is not 0\,\%.} An operator present with 0.3\,\% of a
        filière is present. Those cells read \file{<1\%}.
\end{enumerate}

% ===========================================================================
\section{Known limitations and next steps}

\begin{itemize}
  \item \textbf{No collection dataset.} Every rate on the France page is
        \file{treated / put on market}, a proxy. ADEME publishes \emph{REP ---
        Tonnages issus des lieux de collecte} separately; adding it would make
        the rates real and let the gap-to-target chart use published rates
        rather than a computed proxy.
  \item \textbf{The gap-to-target values are placeholders.} They live in
        \file{lib/data.PLACEHOLDER\_TARGETS} and should be replaced with the
        published figures from the ADEME filière fact sheets --- most of which
        are now in \file{targets\_fr\_br.csv} and could be read from there
        instead.
  \item \textbf{Nothing is municipal.} Per the project's granularity rule, the
        finest available level is \file{dep\_site\_trt}, the department of the
        \emph{treatment site}. Municipal analysis would need a different source.
  \item \textbf{Dark mode is partial.} Charts and cards follow the toggle;
        Streamlit's own widgets do not (Section~\ref{sec:model}).
  \item \textbf{Brazil's side has no equivalent to the French filière data.}
        The targets page compares commitments; there is no Brazilian counterpart
        to \file{REP.csv} in the project yet, so the market-structure analysis
        is France-only.
  \item \textbf{Pre-signed URLs expire.} If the deployed app works today and
        returns 403 next week, that is the first thing to check.
\end{itemize}

\newpage
% ===========================================================================
\appendix
\section{Targets CSV schema}\label{sec:schema}

Plain CSV, comma-separated, UTF-8, one header row, 19 columns in any order. Save
from Excel as ``CSV UTF-8 (Comma delimited)'' --- a French or Portuguese locale
defaults to semicolons, which breaks the read.

\begin{longtable}{@{}rL{3.1cm}L{2.0cm}L{7.0cm}@{}}
\toprule
\textbf{\#} & \textbf{Column} & \textbf{Required} & \textbf{Notes} \\
\midrule
\endhead
1 & \file{country} & yes & \file{France} or \file{Brazil} exactly. Any other
value rejects the whole file \\
2 & \file{level} & yes & \file{government} or \file{sector}, lowercase. Chooses
the page section \\
3 & \file{tier} & no & e.g. \file{EU}, \file{National}, \file{Federal},
\file{Sector agreement} \\
4 & \file{sector} & no & \textbf{The pairing key.} Must match exactly between
the two countries \\
5 & \file{metric\_family} & no & One of the six normalised families \\
6 & \file{metric} & no & The metric as the instrument words it \\
7 & \file{basis} & no & \textbf{The denominator, in words.} Load-bearing \\
8 & \file{unit} & for charts & \file{\%}, \file{Mt}, \file{kg/cap},
\file{count}. Only \file{\%} rows reach the bar charts \\
9--14 & the six numeric fields & no & \file{baseline\_value/year},
\file{latest\_value/year}, \file{target\_value/year} \\
15 & \file{instrument} & no & Cite the instrument, not a summary \\
16 & \file{comparable} & no & \file{TRUE}/\file{FALSE}. Anything unrecognised
is treated as \file{FALSE} \\
17 & \file{status} & yes & \file{placeholder}, \file{sourced} or
\file{verified} \\
18 & \file{note} & no & Caveats, disagreements between sources \\
19 & \file{source\_url} & no & Where the figure came from \\
\bottomrule
\end{longtable}

The six metric families are: Recycling rate; Separate collection rate; Landfill
diversion; Collection coverage; Reuse \& preparation for reuse;
Reverse-logistics recovery. A value outside this list is filtered out of the
page and effectively invisible.

\paragraph{Rejection vs.\ silent coercion.} The loader rejects the whole file
(falling back to placeholders, with the reason shown in red) if a column is
missing or \file{country}/\file{level} contains an unexpected value. It does
\emph{not} reject, but silently blanks, a non-numeric value in a numeric
column --- so \file{50\%}, \file{\textasciitilde 50} and \file{n/a} all quietly
become empty. Write bare numbers.

\section{Troubleshooting}

\begin{longtable}{@{}L{5.4cm}L{9.0cm}@{}}
\toprule
\textbf{Symptom} & \textbf{Cause and fix} \\
\midrule
\endhead
\file{FileNotFoundError}, message redacted, on the deployed app &
The data files are not in the deployment. Either commit \file{data/} or
configure the secrets URLs (Section~\ref{sec:sources}). The pages now catch this
and print which file is missing \\
\addlinespace
Targets page shows the placeholder banner &
No \file{targets\_fr\_br.csv} found locally and no \file{targets\_url} in
secrets \\
\addlinespace
Red ``found but rejected'' banner &
A column is missing, or \file{country}/\file{level} has an unexpected value.
The message names the problem \\
\addlinespace
Charts render white in dark mode &
A \file{st.plotly\_chart} call is missing \file{theme=None}, or the figure does
not set \file{paper\_bgcolor} explicitly \\
\addlinespace
A control appears below the charts it drives &
It was declared after \file{st.columns()}. Move it above the column split \\
\addlinespace
Deployed app returns 403 on data &
A pre-signed URL has expired. Regenerate it and update the secrets \\
\addlinespace
A filière is missing from a chart &
Check the caption underneath --- pages name what they dropped and why \\
\bottomrule
\end{longtable}

\end{document}
